
\documentclass[11pt,a4paper]{article}

\usepackage[margin=2.25cm]{geometry}
\usepackage{amsmath,amssymb,mathtools}
\usepackage{booktabs,longtable,array,tabularx}
\usepackage{enumitem}
\usepackage{hyperref}
\usepackage{xcolor}
\usepackage{listings}
\usepackage{microtype}
\usepackage{fancyhdr}
\usepackage{titlesec}
\usepackage{caption}

\definecolor{navy}{RGB}{24,49,83}
\definecolor{lightgray}{RGB}{245,246,248}
\definecolor{midgray}{RGB}{90,98,108}

\hypersetup{
  colorlinks=true,
  linkcolor=navy,
  urlcolor=navy,
  citecolor=navy,
  pdftitle={Building a First-to-Default Pricer from Bootstrapped CDS Curves},
  pdfauthor={Nigel Li}
}

\pagestyle{fancy}
\fancyhf{}
\lhead{\textcolor{midgray}{FTD Pricer Design Guide}}
\rhead{\textcolor{midgray}{Gaussian Copula}}
\cfoot{\thepage}

\titleformat{\section}{\Large\bfseries\color{navy}}{\thesection}{0.7em}{}
\titleformat{\subsection}{\large\bfseries\color{navy}}{\thesubsection}{0.7em}{}
\titleformat{\subsubsection}{\normalsize\bfseries}{\thesubsubsection}{0.7em}{}

\setlist[itemize]{leftmargin=1.4em,itemsep=0.25em,topsep=0.35em}
\setlist[enumerate]{leftmargin=1.6em,itemsep=0.3em,topsep=0.35em}

\newcommand{\DF}{D(0,t)}
\newcommand{\Q}{\mathbb{Q}}
\newcommand{\E}{\mathbb{E}^{\mathbb{Q}}}
\newcommand{\ind}{\mathbf{1}}
\newcommand{\FTD}{\mathrm{FTD}}
\newcommand{\PV}{\mathrm{PV}}

\lstset{
  basicstyle=\ttfamily\small,
  backgroundcolor=\color{lightgray},
  frame=single,
  rulecolor=\color{lightgray},
  breaklines=true,
  columns=fullflexible,
  showstringspaces=false
}

\begin{document}

\begin{titlepage}
\centering
\vspace*{2.0cm}
{\Huge\bfseries\color{navy} Building a First-to-Default Pricer\par}
\vspace{0.4cm}
{\Large From Bootstrapped CDS Curves to a Gaussian-Copula Basket Model\par}
\vspace{1.3cm}
{\large Implementation guide based on the existing CDS curve and pricer workflow\par}
\vfill
\begin{tabular}{rl}
\textbf{Core inputs:} & SOFR discount factors, single-name CDS curves, recoveries \\
\textbf{Dependence model:} & One-factor or full-matrix Gaussian copula \\
\textbf{Primary engine:} & Monte Carlo, benchmarked by one-factor quadrature \\
\textbf{Target product:} & First-to-default and general $i$-th-to-default baskets
\end{tabular}
\vfill
{\large Nigel Li\par}
{\small July 2026\par}
\end{titlepage}

\tableofcontents
\newpage

\section{Purpose and Scope}

This guide explains how to extend an existing single-name CDS curve bootstrap and CDS pricer into a first-to-default (FTD) basket pricer. The existing CDS work already supplies the marginal default distribution for each reference entity. The additional FTD model must introduce a joint dependence structure, determine the first default time and identity, and value the resulting premium and protection legs.

The implementation should be designed in two layers:

\begin{enumerate}
    \item \textbf{Single-name layer:} discount factors, piecewise hazard rates, survival probabilities, default probabilities, premium-leg PV01 and protection-leg valuation.
    \item \textbf{Basket layer:} correlated default-time generation, first-default selection, basket cash flows, fair FTD spread, correlation calibration and risk.
\end{enumerate}

The intended first production version is a one-factor Gaussian-copula Monte Carlo model with a separate one-dimensional quadrature benchmark. The framework can subsequently be generalised from first-to-default to $i$-th-to-default baskets.

\section{Economic Interpretation of FTD Derivatives}

An FTD basket allows an investor to buy or sell protection on the earliest default in a portfolio of reference names. Since the protection seller is exposed to the first default among several names, the fair FTD spread is normally materially above the simple average single-name spread. Market descriptions sometimes characterise this as approximately two to three times the basket-average spread, but this is only a heuristic. The actual multiple depends on:

\begin{itemize}
    \item the number of names;
    \item each name's CDS term structure;
    \item recovery assumptions;
    \item basket maturity;
    \item correlation;
    \item concentration and heterogeneity;
    \item coupon and accrual conventions.
\end{itemize}

The FTD spread is generally negatively related to default correlation. When correlation is low, each name is a relatively independent opportunity for the basket to experience at least one default. When correlation is high, the names behave more like a single common credit exposure, reducing the probability that at least one name defaults by a given horizon.

For identical names with maturity default probability $p$:

\[
\Q(\tau_{(1)} \leq T)=1-(1-p)^n
\]

under independence, whereas under near-perfect positive dependence the probability approaches $p$ rather than $1-(1-p)^n$.

An $i$-th-to-default basket may contain a small bespoke portfolio or a larger basket, commonly between three and twenty-five names. The FTD is the special case $i=1$.

\section{Existing Single-Name Foundation}

\subsection{Piecewise-Constant Hazard Rates}

For reference name $i$, assume the CDS bootstrap produces the piecewise-constant hazard function

\[
\lambda_i(t)=
\begin{cases}
\lambda_{i,1}, & 0<t\leq T_1,\\
\lambda_{i,2}, & T_1<t\leq T_2,\\
\vdots & \\
\lambda_{i,K}, & T_{K-1}<t\leq T_K.
\end{cases}
\]

For $t\in(T_{k-1},T_k]$, the cumulative hazard is

\[
\Lambda_i(t)
=
\sum_{j=1}^{k-1}
\lambda_{i,j}(T_j-T_{j-1})
+
\lambda_{i,k}(t-T_{k-1}).
\]

The risk-neutral survival probability is

\[
S_i(t)
=
\Q(\tau_i>t)
=
\exp[-\Lambda_i(t)],
\]

and the cumulative default probability is

\[
F_i(t)
=
\Q(\tau_i\leq t)
=
1-S_i(t).
\]

The bucket hazard $\lambda_{i,k}$ is a conditional instantaneous default intensity. Informally, it represents the rate of default during bucket $k$, conditional on survival to that point. It is not itself an unconditional bucket default probability.

The conditional probability of surviving from $T_{k-1}$ to $T_k$ is

\[
\Q(\tau_i>T_k\mid \tau_i>T_{k-1})
=
e^{-\lambda_{i,k}(T_k-T_{k-1})}.
\]

Therefore, the conditional probability of default during that bucket is

\[
1-e^{-\lambda_{i,k}(T_k-T_{k-1})}.
\]

\subsection{What the Existing CDS Pricer Already Provides}

The current CDS curve and pricer framework should already contain, either explicitly or implicitly:

\begin{itemize}
    \item valuation date and settlement conventions;
    \item SOFR discount factors $D(0,t)$;
    \item quoted CDS maturities and par spreads;
    \item recovery assumptions;
    \item premium schedules and day-count fractions;
    \item piecewise hazard rates;
    \item cumulative hazards;
    \item survival and default probabilities;
    \item CDS premium-leg risky PV01;
    \item CDS protection-leg PV;
    \item repricing errors against quoted CDS spreads.
\end{itemize}

These are precisely the marginal inputs required by the Gaussian copula. The FTD extension should reuse them rather than recalibrate the single-name curves inside the basket model.

\section{FTD Contract Definition}

For an $n$-name basket, define each reference entity's default time as $\tau_i$. The first default time is

\[
\tau_{(1)}
=
\min(\tau_1,\ldots,\tau_n).
\]

The identity of the first defaulting name is

\[
K
=
\arg\min_{i\in\{1,\ldots,n\}}\tau_i.
\]

The protection buyer pays the periodic FTD spread $s_{\FTD}$ until the earliest of maturity and first default. If name $K$ defaults first before maturity, the protection buyer receives

\[
N_K(1-R_K),
\]

where $N_K$ is the contractual notional attached to that name and $R_K$ is its recovery rate.

The fair spread solves

\[
\PV_{\mathrm{premium}}(s_{\FTD})
=
\PV_{\mathrm{protection}}.
\]

Since the premium leg is linear in spread,

\[
\boxed{
s_{\FTD}
=
\frac{\PV_{\mathrm{protection}}}
{\mathrm{FTD\ risky\ annuity}}
}
\]

or, in basis points,

\[
s_{\FTD}^{\mathrm{bps}}
=
10{,}000
\frac{\PV_{\mathrm{protection}}}
{\mathrm{FTD\ risky\ annuity}}.
\]

\section{Why a Copula Is Required}

The bootstrapped CDS curves determine each marginal distribution $F_i(t)$, but they do not determine the joint law

\[
\Q(\tau_1\leq t_1,\ldots,\tau_n\leq t_n).
\]

A copula joins the marginal distributions into a joint distribution while preserving each marginal curve. Under a Gaussian copula,

\[
C_{\Sigma}(u_1,\ldots,u_n)
=
\Phi_{\Sigma}
\left(
\Phi^{-1}(u_1),\ldots,\Phi^{-1}(u_n)
\right),
\]

where $\Phi_{\Sigma}$ is the multivariate standard normal CDF with correlation matrix $\Sigma$.

The implementation therefore follows the chain

\[
\text{CDS spreads}
\rightarrow
\lambda_i(t)
\rightarrow
F_i(t)
\rightarrow
\text{Gaussian copula}
\rightarrow
(\tau_1,\ldots,\tau_n)
\rightarrow
\tau_{(1)}
\rightarrow
\text{FTD cash flows}.
\]

\section{Gaussian-Copula Default-Time Simulation}

\subsection{Full Correlation-Matrix Form}

Let

\[
\mathbf Z=(Z_1,\ldots,Z_n)^\top
\sim N(\mathbf 0,\Sigma).
\]

Compute a Cholesky factorisation

\[
\Sigma=LL^\top,
\]

generate independent standard normals

\[
\boldsymbol{\varepsilon}\sim N(\mathbf 0,I),
\]

and set

\[
\mathbf Z=L\boldsymbol{\varepsilon}.
\]

Convert each latent normal to a uniform:

\[
U_i=\Phi(Z_i).
\]

Then convert the uniform into a default time:

\[
\tau_i=F_i^{-1}(U_i).
\]

This construction preserves the calibrated marginal default distribution:

\[
\Q(\tau_i\leq t)=F_i(t).
\]

\subsection{One-Factor Form}

A simpler and highly practical first implementation is

\[
Z_i
=
\beta_iY
+
\sqrt{1-\beta_i^2}\,\varepsilon_i,
\]

where $Y$ and the $\varepsilon_i$ are independent standard normal variables. Then

\[
\operatorname{Corr}(Z_i,Z_j)=\beta_i\beta_j.
\]

In the homogeneous case,

\[
\beta_i=\sqrt{\rho}
\]

for all $i$, giving

\[
Z_i
=
\sqrt{\rho}\,Y
+
\sqrt{1-\rho}\,\varepsilon_i
\]

and pairwise latent correlation $\rho$.

The parameter $\rho$ is a latent Gaussian correlation. It is not directly a CDS-spread correlation, hazard-rate correlation, default-time Pearson correlation or observed equity-return correlation.

\subsection{Default-Time Inversion for Piecewise Hazards}

Given $U_i$, define the exponential threshold

\[
H_i=-\ln(1-U_i).
\]

Since

\[
F_i(t)=1-e^{-\Lambda_i(t)},
\]

the simulated default time solves

\[
\Lambda_i(\tau_i)=H_i.
\]

Find bucket $k$ such that

\[
\Lambda_i(T_{k-1})<H_i\leq\Lambda_i(T_k).
\]

Then

\[
\boxed{
\tau_i
=
T_{k-1}
+
\frac{
H_i-\Lambda_i(T_{k-1})
}{
\lambda_{i,k}
}
}
\]

provided $\lambda_{i,k}>0$.

If $H_i$ lies beyond the final calibrated maturity, the model requires an explicit extrapolation rule. The simplest is a flat terminal hazard:

\[
\lambda_i(t)=\lambda_{i,K},
\qquad t>T_K.
\]

\section{Monte Carlo Valuation}

For path $m$, generate all default times and identify

\[
\tau_{(1)}^{(m)}
=
\min_i\tau_i^{(m)},
\qquad
K^{(m)}
=
\arg\min_i\tau_i^{(m)}.
\]

\subsection{Protection Leg}

The pathwise protection PV is

\[
\PV_{\mathrm{prot}}^{(m)}
=
\ind_{\{\tau_{(1)}^{(m)}\leq T\}}
D(0,\tau_{(1)}^{(m)})
N_{K^{(m)}}(1-R_{K^{(m)}}).
\]

The Monte Carlo estimate is

\[
\PV_{\mathrm{prot}}
\approx
\frac{1}{M}
\sum_{m=1}^{M}
\PV_{\mathrm{prot}}^{(m)}.
\]

\subsection{Scheduled Premium Leg}

For premium dates $t_1,\ldots,t_J$ and accrual fractions $\alpha_j$, the scheduled risky annuity on path $m$ is

\[
A_{\mathrm{sched}}^{(m)}
=
\sum_{j=1}^{J}
\alpha_jD(0,t_j)
\ind_{\{\tau_{(1)}^{(m)}>t_j\}}.
\]

\subsection{Accrued Premium on Default}

If

\[
t_{j-1}<\tau_{(1)}^{(m)}\leq t_j,
\]

the accrued contribution is

\[
A_{\mathrm{accr}}^{(m)}
=
D(0,\tau_{(1)}^{(m)})
\operatorname{DCF}(t_{j-1},\tau_{(1)}^{(m)}).
\]

Therefore,

\[
A_{\FTD}
\approx
\frac{1}{M}
\sum_{m=1}^{M}
\left(
A_{\mathrm{sched}}^{(m)}
+
A_{\mathrm{accr}}^{(m)}
\right).
\]

The fair spread is then

\[
s_{\FTD}
=
\frac{\PV_{\mathrm{prot}}}{A_{\FTD}}.
\]

\section{One-Factor Quadrature Benchmark}

Monte Carlo should be the general engine, but the one-factor model also provides a strong independent benchmark.

Define the default threshold

\[
a_i(t)=\Phi^{-1}(F_i(t)).
\]

Conditional on $Y=y$,

\[
Z_i\mid Y=y
\sim
N(\beta_i y,1-\beta_i^2).
\]

Hence the conditional default probability is

\[
p_i(t\mid y)
=
\Phi\left(
\frac{a_i(t)-\beta_i y}
{\sqrt{1-\beta_i^2}}
\right),
\]

and conditional survival is

\[
q_i(t\mid y)=1-p_i(t\mid y).
\]

Conditional independence gives

\[
\Q(\tau_{(1)}>t\mid Y=y)
=
\prod_{i=1}^{n}q_i(t\mid y).
\]

Integrating over $Y$ gives the basket survival probability

\[
\boxed{
S_{\FTD}(t)
=
\int_{-\infty}^{\infty}
\phi(y)
\prod_{i=1}^{n}
\Phi\left(
\frac{\beta_i y-\Phi^{-1}(F_i(t))}
{\sqrt{1-\beta_i^2}}
\right)
dy
}.
\]

This is a one-dimensional integral regardless of basket size and can be evaluated efficiently with Gauss--Hermite quadrature.

For common recovery and notional,

\[
\PV_{\mathrm{prot}}
=
N(1-R)
\int_0^T D(0,t)\,d[1-S_{\FTD}(t)].
\]

The scheduled annuity is

\[
A_{\mathrm{sched}}
=
\sum_{j=1}^{J}
\alpha_jD(0,t_j)S_{\FTD}(t_j).
\]

A simple midpoint approximation for accrual is

\[
A_{\mathrm{accr}}
\approx
\sum_{j=1}^{J}
D(0,\bar t_j)
\frac{\alpha_j}{2}
\left[
S_{\FTD}(t_{j-1})-S_{\FTD}(t_j)
\right].
\]

\section{Workbook and Sheet Architecture}

\subsection{Important Repository Limitation}

The public repository confirms the broader curve-bootstrapping project and its SOFR/Excel implementation documentation. However, the binary CDS workbook's internal tabs were not available through the repository connector used for this review. Therefore, the tab names below are a professional target architecture and mapping for the CDS Pricer 2.2 workflow, not a claim that these are the exact existing sheet names.

Where an existing sheet already performs the stated function, retain it and add the specified output columns rather than duplicating it.

\subsection{Recommended Existing-Sheet Mapping}

\begin{longtable}{p{3.3cm}p{5.5cm}p{6.0cm}}
\toprule
\textbf{Sheet / role} & \textbf{Current CDS function} & \textbf{FTD use and required extension} \\
\midrule
\endhead

\texttt{Control} &
Valuation date, trade date, settlement date, day-count assumptions, recovery, curve selection and global switches. &
Add basket maturity, basket notional, number of names, correlation model, scalar $\rho$ or matrix selector, simulation count, seed, antithetic switch and extrapolation rule. \\

\texttt{Discount Curve} &
Stores SOFR pillar dates, discount factors, zero rates and interpolation inputs. &
No conceptual change. Expose a clean function or lookup range returning $D(0,t)$ for arbitrary coupon and simulated default dates. \\

\texttt{CDS Market Data} &
Contains single-name CDS spreads by maturity, recovery assumptions and possibly upfront quotes. &
Convert to a basket input table with one row per name and one column per quoted maturity. Add name-specific notional, recovery and sector metadata. \\

\texttt{CDS Schedule} &
Builds coupon dates, accrual fractions, protection start and maturity conventions. &
Create one master FTD premium schedule because the basket terminates on first default. Confirm all names use compatible contract dates or define a common basket schedule explicitly. \\

\texttt{Hazard Bootstrap} &
Solves sequentially for piecewise hazards that reprice each CDS quote. &
Run independently for every basket name. Output a rectangular table indexed by name and hazard bucket. \\

\texttt{Survival Curve} &
Calculates cumulative hazard, survival probability and default probability at curve nodes or coupon dates. &
Expose $F_i(t)$ and $F_i^{-1}(u)$. Add cumulative-hazard breakpoint arrays used by default-time inversion. \\

\texttt{CDS Pricer} &
Calculates premium leg, accrued premium, protection leg, par spread and PV for one name. &
Use as a single-name validation engine. A one-name FTD must reproduce this sheet's par CDS spread under identical conventions. \\

\texttt{Calibration Check} &
Shows quote repricing errors and curve monotonicity. &
Add basket-wide pass/fail flags: non-negative hazards, decreasing survival, marginal simulation checks and maximum repricing error. \\

\bottomrule
\end{longtable}

\subsection{New FTD Sheets to Add}

\begin{longtable}{p{3.4cm}p{5.4cm}p{6.0cm}}
\toprule
\textbf{New sheet} & \textbf{Purpose} & \textbf{Principal contents} \\
\midrule
\endhead

\texttt{FTD Basket Input} &
Defines the product and constituent-level terms. &
Name, weight/notional, recovery, CDS curve reference, maturity, currency, seniority, correlation loading $\beta_i$, inclusion flag. \\

\texttt{FTD Correlation} &
Stores and validates the dependence model. &
Scalar $\rho$, one-factor loadings, or full matrix $\Sigma$; symmetry checks; diagonal checks; eigenvalues; Cholesky factor. \\

\texttt{FTD Randoms} &
Generates systemic and idiosyncratic random variables. &
$Y$, $\varepsilon_i$, $Z_i$, $U_i$ by simulation path. Prefer an external Python engine for large path counts. \\

\texttt{FTD Default Times} &
Maps copula uniforms through each marginal curve. &
$H_i=-\ln(1-U_i)$, selected hazard bucket, inverted $\tau_i$, first default time and first default identity. \\

\texttt{FTD Cash Flows} &
Calculates pathwise premium and protection legs. &
Scheduled survival indicators, accrued premium, loss-given-default, discount factors and pathwise PVs. \\

\texttt{FTD Results} &
Reports price and fair spread. &
Protection PV, risky annuity, fair spread, standard error, confidence interval, expected first-default time and first-name contribution. \\

\texttt{FTD Sensitivities} &
Calculates risk and scenario analysis. &
Correlation01, constituent CS01, recovery sensitivity, IR01, basket-name contribution and correlation grid. \\

\texttt{FTD Validation} &
Runs theoretical and numerical test cases. &
One-name equivalence, zero-correlation identity, near-one-correlation limit, marginal preservation and Monte Carlo versus quadrature. \\

\bottomrule
\end{longtable}

\section{Detailed Excel Implementation}

\subsection{Step 1: Standardise the Name-by-Bucket Hazard Table}

Create a table with one row per name and two aligned blocks:

\begin{enumerate}
    \item hazard rates $\lambda_{i,k}$;
    \item cumulative hazards $\Lambda_i(T_k)$.
\end{enumerate}

For name $i$ and bucket $k$:

\[
\Lambda_i(T_k)
=
\Lambda_i(T_{k-1})
+
\lambda_{i,k}(T_k-T_{k-1}).
\]

In Excel, conceptually:

\begin{lstlisting}
=CumulativeHazard_Previous
 + Hazard_Current * (BucketEnd - BucketStart)
\end{lstlisting}

Ensure all time differences use the same year-fraction convention as the bootstrap.

\subsection{Step 2: Add Inverse Default-Time Logic}

For each path and name:

\begin{lstlisting}
U = NORM.S.DIST(Z, TRUE)
H = -LN(1-U)
\end{lstlisting}

Use \texttt{XMATCH} or \texttt{MATCH} to find the first cumulative-hazard breakpoint greater than or equal to $H$.

Conceptually:

\begin{lstlisting}
Bucket = XMATCH(H, CumulativeHazardRow, 1)
Tau = BucketStart
    + (H - CumHazardAtBucketStart) / HazardInBucket
\end{lstlisting}

Care is required because approximate-match modes differ across Excel versions. Test boundary values exactly equal to each cumulative-hazard breakpoint.

\subsection{Step 3: Generate Correlated Normals}

For a homogeneous one-factor model:

\begin{lstlisting}
SystemicY = NORM.S.INV(RAND())
IdioEps  = NORM.S.INV(RAND())
Z_i      = SQRT(rho)*SystemicY + SQRT(1-rho)*IdioEps
\end{lstlisting}

For reproducibility, ordinary volatile \texttt{RAND()} formulas are not ideal. A Python engine should generate and save a fixed random-number matrix using a specified seed.

\subsection{Step 4: Select the First Default}

For each path:

\begin{lstlisting}
FirstDefaultTime = MIN(DefaultTime_Name1:DefaultTime_NameN)
FirstNameIndex   = XMATCH(FirstDefaultTime,
                          DefaultTime_Name1:DefaultTime_NameN,
                          0)
\end{lstlisting}

Simultaneous defaults have zero probability under a continuous Gaussian-copula model. Nevertheless, implement a deterministic tie rule for numerical robustness.

\subsection{Step 5: Calculate Pathwise Protection PV}

\begin{lstlisting}
=IF(FirstDefaultTime<=Maturity,
    DF_at_FirstDefault
    * Notional_of_FirstName
    * (1-Recovery_of_FirstName),
    0)
\end{lstlisting}

The discount factor at a simulated date must be interpolated from the SOFR curve using the same interpolation convention as the existing pricer.

\subsection{Step 6: Calculate Scheduled Premium PV01}

For coupon date $t_j$:

\begin{lstlisting}
=AccrualFactor_j
 * DF_j
 * --(FirstDefaultTime > CouponDate_j)
\end{lstlisting}

Sum across all coupon dates.

\subsection{Step 7: Calculate Accrued Premium}

Find the coupon interval containing the first default time. Then calculate

\begin{lstlisting}
=IF(FirstDefaultTime<=Maturity,
    DF_at_FirstDefault
    * YearFrac(PreviousCouponDate, FirstDefaultTime),
    0)
\end{lstlisting}

Use the contractual accrual convention, not a simple calendar-day fraction unless that is exactly the convention in the current CDS pricer.

\subsection{Step 8: Aggregate and Solve Fair Spread}

\begin{lstlisting}
ProtectionPV = AVERAGE(PathwiseProtectionPV)
RiskyAnnuity = AVERAGE(PathwiseScheduledPV01
                     + PathwiseAccruedPV01)
FairSpread   = ProtectionPV / RiskyAnnuity
FairSpreadBp = 10000 * FairSpread
\end{lstlisting}

Also report the Monte Carlo standard error:

\[
\operatorname{SE}(\bar X)
=
\frac{s_X}{\sqrt{M}}.
\]

For the ratio defining the fair spread, use either common-path bumping or a delta-method estimate if a spread confidence interval is required.

\section{Recommended Python Calculation Engine}

Excel is appropriate for transparency, controls and reporting. The simulation itself should preferably be implemented in Python and returned to Excel as summarised outputs.

\subsection{Core Data Structures}

\begin{lstlisting}[language=Python]
@dataclass
class HazardCurve:
    breakpoints: np.ndarray
    hazards: np.ndarray
    cumulative_hazards: np.ndarray
    recovery: float

@dataclass
class FTDBasket:
    names: list[str]
    notionals: np.ndarray
    curves: list[HazardCurve]
    maturity: float
    premium_dates: np.ndarray
    accrual_fractions: np.ndarray
\end{lstlisting}

\subsection{Default-Time Inversion}

\begin{lstlisting}[language=Python]
def inverse_default_times(u, breakpoints, hazards, cumulative):
    h = -np.log1p(-u)
    k = np.searchsorted(cumulative, h, side="left")

    tau = np.empty_like(h)

    inside = k < len(hazards)
    k_in = k[inside]

    previous_cum = np.where(
        k_in == 0, 0.0, cumulative[k_in - 1]
    )
    previous_time = np.where(
        k_in == 0, 0.0, breakpoints[k_in - 1]
    )

    tau[inside] = (
        previous_time
        + (h[inside] - previous_cum) / hazards[k_in]
    )

    beyond = ~inside
    tau[beyond] = (
        breakpoints[-1]
        + (h[beyond] - cumulative[-1]) / hazards[-1]
    )

    return tau
\end{lstlisting}

\subsection{One-Factor Simulation}

\begin{lstlisting}[language=Python]
rng = np.random.default_rng(seed)

y = rng.standard_normal((n_paths, 1))
eps = rng.standard_normal((n_paths, n_names))

z = beta.reshape(1, -1) * y + np.sqrt(
    1.0 - beta.reshape(1, -1) ** 2
) * eps

u = scipy.stats.norm.cdf(z)
\end{lstlisting}

Use vectorised inversion by name, obtain the row-wise minimum, and calculate pathwise cash flows.

\section{Correlation Input and Calibration}

Single-name CDS quotes do not determine correlation. Correlation must come from an additional market input or scenario assumption.

\subsection{Implied FTD Correlation}

If a market FTD spread is available, solve

\[
s_{\mathrm{model}}(\rho)
-
s_{\mathrm{market}}
=
0.
\]

Use bisection or Brent's method on a bounded interval such as

\[
\rho\in[0,0.999].
\]

The model spread should generally decline with $\rho$ for an FTD basket, but verify monotonicity numerically for heterogeneous portfolios and contractual features.

\subsection{Alternative Inputs}

Possible alternatives include:

\begin{itemize}
    \item historical equity or asset-return correlations;
    \item sector-factor models;
    \item index-tranche-implied dependence;
    \item a conservative scenario grid.
\end{itemize}

Historical correlations are physical-measure estimates and are not automatically risk-neutral. They should be labelled as assumptions rather than market calibrations.

\section{Risk Measures}

The minimum risk report should include:

\begin{description}
    \item[Correlation01:] Change in PV or fair spread for a one-percentage-point change in $\rho$.
    \item[Constituent CS01:] Change in basket PV after bumping one name's CDS spread curve by one basis point and re-bootstrapping that name.
    \item[Parallel basket CS01:] Simultaneous one-basis-point bump to all constituent CDS curves.
    \item[Recovery sensitivity:] Change in PV after changing one or all recovery assumptions.
    \item[IR01:] Change in PV after a one-basis-point shift to the discount curve.
    \item[Name contribution:] Expected discounted loss contribution by first-defaulting name.
\end{description}

Use common random numbers for all bumped Monte Carlo calculations. This materially reduces noise in finite-difference sensitivities.

\section{Validation Framework}

\subsection{Single-Name Equivalence}

For $n=1$, the FTD fair spread must equal the existing single-name CDS par spread, subject to identical timing, accrual, discounting and recovery conventions.

\subsection{Zero-Correlation Check}

At $\rho=0$,

\[
S_{\FTD}(t)
=
\prod_{i=1}^{n}S_i(t).
\]

The Monte Carlo result should converge to this analytical identity.

\subsection{Perfect-Dependence Limit}

For identical names and $\rho\rightarrow1$, the FTD should approach a single-name CDS because all latent variables and default times converge.

\subsection{Marginal Preservation}

For each name and date $t$,

\[
\frac{1}{M}
\sum_{m=1}^{M}
\ind_{\{\tau_i^{(m)}\leq t\}}
\approx
F_i(t).
\]

Changing correlation must not alter the unconditional marginal default curve.

\subsection{Monte Carlo versus Quadrature}

Under the one-factor model with common recovery and notional, Monte Carlo and Gauss--Hermite quadrature should agree within Monte Carlo sampling error.

\subsection{Correlation-Matrix Checks}

For a full matrix model, require:

\[
\Sigma=\Sigma^\top,
\qquad
\Sigma_{ii}=1,
\qquad
\Sigma\succeq0.
\]

If the matrix is not positive semidefinite, Cholesky decomposition will fail or produce an invalid dependence structure.

\section{Extension to $i$-th-to-Default}

The simulation framework generalises directly. Sort each path's default times:

\[
\tau_{(1)}
\leq
\tau_{(2)}
\leq
\cdots
\leq
\tau_{(n)}.
\]

For an $i$-th-to-default contract, termination and protection payment occur at $\tau_{(i)}$ rather than $\tau_{(1)}$.

In Python:

\begin{lstlisting}[language=Python]
ith_time = np.partition(default_times, i - 1, axis=1)[:, i - 1]
ith_name = np.argpartition(default_times, i - 1, axis=1)[:, i - 1]
\end{lstlisting}

The premium leg survives until the $i$-th default. The protection payment uses the recovery and notional of the $i$-th defaulting name. For $i>1$, correlation sensitivity can differ substantially from the FTD case because the product requires multiple defaults rather than merely one.

\section{Recommended Build Order}

\begin{enumerate}
    \item Freeze and validate the existing SOFR discount curve.
    \item Validate each single-name hazard curve by repricing its quoted CDS spreads.
    \item Refactor the hazard output into a standard name-by-bucket table.
    \item Implement and unit-test inverse default-time mapping.
    \item Build a homogeneous one-factor Gaussian-copula simulator.
    \item Add first-default time and identity selection.
    \item Add scheduled premium, accrued premium and protection cash flows.
    \item Produce fair spread, PV and Monte Carlo error.
    \item Implement the one-factor quadrature benchmark.
    \item Add correlation scenarios and implied-correlation calibration.
    \item Add constituent CS01, correlation01, recovery risk and IR01.
    \item Generalise from first-to-default to $i$-th-to-default.
\end{enumerate}

\section{Model Limitations}

The Gaussian copula is useful because it is tractable and preserves calibrated single-name marginals. Its limitations must remain visible:

\begin{itemize}
    \item a constant latent correlation does not capture correlation skew;
    \item the Gaussian copula has no lower-tail dependence beyond that implied by linear Gaussian correlation;
    \item static default-time copulas do not model spread dynamics;
    \item contagion after a realised default is not represented;
    \item historical correlations may differ materially from risk-neutral dependence;
    \item recovery is often treated as deterministic even though recovery and default dependence may be material;
    \item bespoke-basket correlation calibration may be underdetermined.
\end{itemize}

The model should therefore report correlation scenarios and concentration measures rather than presenting one correlation assumption as uniquely correct.

\section{Final Target Architecture}

The completed system should have the following separation:

\[
\boxed{
\begin{array}{c}
\text{SOFR discount curve} \\
+ \\
\text{single-name CDS bootstraps} \\
\downarrow \\
\text{validated marginal survival curves} \\
\downarrow \\
\text{Gaussian dependence engine} \\
\downarrow \\
\text{correlated default times} \\
\downarrow \\
\text{first or $i$-th default event} \\
\downarrow \\
\text{premium and protection cash flows} \\
\downarrow \\
\text{PV, fair spread, calibration and risk}
\end{array}
}
\]

The existing CDS Pricer 2.2 should remain the authoritative single-name calibration and validation layer. The FTD module should consume its hazard and survival outputs, not recreate them independently. This gives a transparent lineage from market CDS spreads to each constituent's marginal curve and then to the basket valuation.

\section*{References}

\begin{enumerate}
    \item D. X. Li, ``On Default Correlation: A Copula Function Approach,'' \textit{Journal of Fixed Income}, 2000. Available through SSRN: \url{https://papers.ssrn.com/sol3/papers.cfm?abstract_id=187289}.
    \item ISDA CDS Standard Model, documentation and production model resources: \url{https://www.cdsmodel.com/documentation.html}.
    \item N. Li, \texttt{curve-bootstrapping} repository: \url{https://github.com/nl2992/curve-bootstrapping}.
\end{enumerate}

\end{document}
